\chapter{Obligations d'État}\label{ch:m2:government-bonds}

Deux traders, sur un salon de discussion, conviennent d'échanger 10~millions de dollars d'un \omterm{def:m2:money-markets:bill}{bon du Trésor} à dix ans à un
rendement de 4.200~\%. Quand les confirmations reviennent, elles diffèrent de 47\,351~dollars. Les deux ont
évalué le même bon au même rendement, à la même date de règlement-livraison, et tous deux
ont raison~: l'un a retenu le cours affiché à l'écran, 100-12+, et l'autre
le montant qui change de mains, auquel s'ajoute le \omterm{def:m2:government-bonds:accrued}{coupon couru} depuis le
dernier paiement, soit quarante et un jours d'un \omterm{def:m2:government-bonds:coupon}{coupon} de 4.25~\%. Le bon est l'instrument le plus
simple de ce livre~: un échéancier fixe de paiements versés par un État qui
emprunte dans sa propre monnaie. Tout ce qui le concerne, son cours affiché, son
rendement, sa sensibilité aux taux, dépend de conventions fixées par
la réglementation et les usages de marché, et que tout système qui y touche doit
les appliquer sans erreur. Ce chapitre les expose et dérive les mesures de risque sur
lesquelles repose la suite de ce livre.

\section{Flux de trésorerie, conventions de décompte des jours et coupon couru}

\begin{definition}[Obligation à coupon]\label{def:m2:government-bonds:coupon}
Une obligation à coupon fixe verse, pour 100 de valeur nominale, un \emph{coupon}\index{coupon}
de $c/f$ à chacune des $f$ dates par an et rembourse sa valeur nominale, appelée
\emph{pair}\index{pair}, à l'échéance. Les dates de coupon sont engendrées à rebours
à partir de la date d'échéance, par pas de $12/f$ mois~; si l'échéance est le dernier
jour d'un mois, il en est de même de chaque date de coupon.
\end{definition}

\begin{definition}[Convention de décompte des jours]\label{def:m2:government-bonds:daycount}
Une \emph{convention de décompte des jours}\index{convention de décompte des jours} transforme deux dates en
une fraction d'année. \emph{Actual/360} divise le nombre réel de jours par
360~; \emph{actual/365 fixed} par 365~; \emph{30/360} compte chaque mois pour 30
jours~; \emph{actual/actual} (ICMA) compte les jours réels écoulés dans une période de \omterm{def:m2:government-bonds:coupon}{coupon}
et divise par le nombre de jours réels de cette période, de sorte que chaque période complète
rapporte exactement $c/f$.
\end{definition}

\begin{definition}[Coupon couru, prix pied de coupon et prix plein coupon]\label{def:m2:government-bonds:accrued}
Entre deux dates de \omterm{def:m2:government-bonds:coupon}{coupon}, le vendeur d'une obligation a acquis une partie du
\omterm{def:m2:government-bonds:coupon}{coupon} suivant, et l'acheteur le lui verse~: c'est le \emph{coupon
couru}\index{coupon couru}, $A = \frac{c}{f}\,(1 - w)$ sous
actual/actual, où $w$ est la fraction de la période courante restant à courir.
Le \emph{prix plein coupon}\index{prix plein coupon} (ou prix de facturation) est ce qui
change de mains~; le \emph{prix pied de coupon}\index{prix pied de coupon} affiché à l'écran est
le prix plein \omterm{def:m2:government-bonds:coupon}{coupon} diminué du \omterm{def:m2:government-bonds:coupon}{coupon} couru.
\end{definition}

Le marché cote les \omterm{def:m2:government-bonds:accrued}{prix pied de coupon} parce qu'ils ne sautent pas aux dates de \omterm{def:m2:government-bonds:coupon}{coupon}
(voir la \cref{fig:m2:government-bonds:cleandirty})~: un \omterm{def:m2:government-bonds:accrued}{prix pied de coupon} qui bouge signifie
que le marché a bougé. Sur le marché américain, les \omterm{def:m2:government-bonds:accrued}{prix pied de coupon} des \omterm{def:m2:money-markets:bill}{bons du Trésor} et des obligations
sont cotés en trente-deuxièmes de point, de sorte que 100-12 vaut $100 + 12/32 =
100.375$, et un signe plus final ajoute un soixante-quatrième~: 100-12+ vaut 100.390625.

\begin{example}[Quarante et un jours]\label{ex:m2:government-bonds:accrued}
Un \omterm{def:m2:money-markets:bill}{bon du Trésor} versant un \omterm{def:m2:government-bonds:coupon}{coupon} de 4.25~\% par semestre les 15 février et 15 août est réglé le 25
septembre 2026. La période courante compte 184 jours, dont 41 sont écoulés~:
$A = 2.125 \times 41/184 = 0.473505$ pour 100, soit 4\,735~dollars par million, 47\,351~dollars
sur la transaction du paragraphe d'ouverture.
\end{example}

\begin{omfigure}
\begin{tikzpicture}[font=\small]
\draw[thick,-{Stealth[length=2.5mm]}] (0,0) -- (13.2,0) node[right]{temps};
\foreach \x/\t/\s in {1.0/coupon précédent/15 août 2026,6.0/coupon suivant/15 févr. 2027,11.0/coupon d'après/15 août 2027}{
  \draw[thick] (\x,0.15) -- (\x,-0.15);
  \node[below=4pt,font=\footnotesize,align=center] at (\x,-0.15) {\t\\\s};}
\draw[omThm,very thick] (2.25,0.25) -- (2.25,-0.25);
\draw[omThm,densely dotted] (2.25,0.25) -- (2.25,1.1);
\node[anchor=south,font=\footnotesize\bfseries,text=omThm] at (2.25,1.1) {règlement, 25 sept.};
\draw[decorate,decoration={brace,amplitude=4pt},omProp,thick] (1.0,0.3) -- (2.25,0.3) node[midway,above=5pt,font=\footnotesize,text=omProp] {$1 - w$};
\draw[decorate,decoration={brace,amplitude=4pt},omDef,thick] (2.25,0.3) -- (6.0,0.3) node[midway,above=5pt,font=\footnotesize,text=omDef] {$w = 143/184$};
\foreach \x in {6.0,11.0}{\draw[omDef,-{Stealth[length=2mm]},thick] (\x,0.6) -- (\x,1.4);}
\node[font=\footnotesize,text=omDef,align=center] at (8.5,1.35) {coupons de $c/f$, actualisés\\sur $w$, $1 + w$, \dots\ périodes};
\end{tikzpicture}

\omcaption{Un règlement entre deux dates de \omterm{def:m2:government-bonds:coupon}{coupon}. Le vendeur reçoit la
fraction $1 - w$ du \omterm{def:m2:government-bonds:coupon}{coupon} à venir au titre du \omterm{def:m2:government-bonds:accrued}{coupon couru}~; l'acheteur reçoit
le \omterm{def:m2:government-bonds:coupon}{coupon} entier. La convention de marché actualise chaque flux ultérieur sur $w$,
$1 + w$, $2 + w$, \dots\ périodes.}\label{fig:m2:government-bonds:timeline}
\end{omfigure}

\begin{dated}{2026-09}{Conventions de trois marchés souverains}\label{dat:m2:government-bonds:conventions}
\textbf{Bons et obligations du Trésor américain}~: \omterm{def:m2:government-bonds:coupon}{coupons} semestriels aux dates fixées par
l'échéance~; chaque paiement régulier représente exactement la moitié du \omterm{def:m2:government-bonds:coupon}{coupon} annuel~; le \omterm{def:m2:government-bonds:accrued}{coupon
couru} se calcule sur les jours réels du semestre~; les prix sont cotés en trente-deuxièmes de point.
\textbf{Gilts britanniques}~: \omterm{def:m2:government-bonds:coupon}{coupons} semestriels~; \omterm{def:m2:government-bonds:accrued}{coupon couru} actual/actual depuis
novembre 1998~; détachement du dividende sept jours ouvrés avant chaque paiement, après
quoi l'acheteur ne perçoit pas le \omterm{def:m2:government-bonds:coupon}{coupon} à venir. \textbf{Bunds allemands}~:
émis à 7, 10, 15 ou 30 ans, avec un \omterm{def:m2:government-bonds:coupon}{coupon} \emph{annuel} fixe.
\end{dated}

\section{Prix et rendement}

\begin{definition}[Rendement actuariel]\label{def:m2:government-bonds:ytm}
Le \emph{rendement actuariel}\index{rendement actuariel} $y$ d'une obligation comptant $n$
\omterm{def:m2:government-bonds:coupon}{coupons} restants, réglée à la fraction $w$ d'une période avant le suivant, est le taux
capitalisé $f$ fois par an, qui égalise le \omterm{def:m2:government-bonds:accrued}{prix plein coupon} avec
la somme de ses flux actualisés~:
\[
P_{\text{plein}} \;=\; \sum_{k=0}^{n-1} \frac{c/f}{(1 + y/f)^{k+w}} \;+\; \frac{100}{(1 + y/f)^{n-1+w}} .
\]
\end{definition}

\begin{proposition}[Prix et rendement]\label{prop:m2:government-bonds:priceyield}
Le \omterm{def:m2:government-bonds:accrued}{prix plein coupon} est une fonction strictement décroissante et strictement convexe de $y$. À
une date de \omterm{def:m2:government-bonds:coupon}{coupon} ($w = 1$ pour la période qui débute), le prix vaut 100 si et
seulement si $y = c/100$.
\end{proposition}
\begin{proof}
Chaque terme $a_k (1 + y/f)^{-t_k}$ avec $a_k, t_k > 0$ est strictement décroissant
et strictement convexe en $y$, il en est de même de leur somme. Avec $w = 1$ et $v = 1/(1 +
y/f)$, le prix vaut $\frac cf \sum_{k=1}^{n} v^k + 100 v^n$~; avec $y/f =
c/(100 f)$, l'annuité vaut $(1 - v^n)/(y/f)$ et le prix vaut $100(1 - v^n)
+ 100 v^n = 100$. L'unicité découle de la monotonie.
\end{proof}

\begin{remark}[Deux conventions pour la fraction]\label{rem:m2:government-bonds:conventions}
La formule ci-dessus, qui capitalise sur la fraction $w$, est celle du marché.
Le règlement propre du Trésor américain, utilisé pour fixer les prix d'adjudication, actualise la
période fractionnaire en intérêts simples, $P_{\text{plein}}(1 + w\,y/f) =$ valeur à
la date du \omterm{def:m2:government-bonds:coupon}{coupon} suivant. Les deux coïncident aux dates de \omterm{def:m2:government-bonds:coupon}{coupon} et diffèrent
de quelques millièmes de point entre elles. Un système qui reproduit les
exemples travaillés du règlement et les écrans du marché doit implémenter les deux
(\cref{bld:m2:government-bonds:bond}).
\end{remark}

\begin{dated}{2026-09}{Dette négociable américaine}\label{dat:m2:government-bonds:debt}
Encours au 31 août 2026~: 31.8~billions de dollars de titres du Trésor
négociables, dont 7.2 billions de \omterm{def:m2:money-markets:bill}{bons du Trésor}, 16.2 billions de billets, 5.5
billions d'obligations, 2.2 billions de titres indexés sur l'inflation et des billets
à taux variable, 0.7~billion.
\end{dated}

\section{Duration, DV01 et convexité}

\begin{definition}[Duration, DV01, convexité]\label{def:m2:government-bonds:duration}
Avec des flux $a_k$ aux dates $t_k$ (en périodes) et $v = 1/(1 + y/f)$, la
\emph{duration de Macaulay}\index{duration de Macaulay} est la moyenne des dates
des flux pondérée par leur valeur actuelle, $D = \frac{1}{f P}\sum_k t_k a_k v^{t_k}$,
en années. La \emph{duration modifiée}\index{duration modifiée} vaut
$D_{\mathrm{mod}} = -\frac1P \frac{dP}{dy} = D/(1 + y/f)$.
Le \emph{DV01}\index{DV01} est la baisse du prix pour une hausse d'un point de base
du rendement, $P\,D_{\mathrm{mod}} \times 10^{-4}$ pour 100, généralement coté par million
de valeur nominale. La \emph{convexité}\index{convexité} vaut $\mathcal C = \frac1P
\frac{d^2P}{dy^2}$.
\end{definition}

\begin{proposition}[Variation de prix au second ordre]\label{prop:m2:government-bonds:taylor}
Pour une variation parallèle $\Delta y$ du rendement,
\[
\frac{\Delta P}{P} \;=\; -D_{\mathrm{mod}}\,\Delta y \;+\; \tfrac12\,\mathcal C\,(\Delta y)^2 \;+\; O(\Delta y^3),
\qquad
\mathcal C \;=\; \frac{1}{f^2 P (1+y/f)^2}\sum_k t_k(t_k+1)\,a_k v^{t_k} .
\]
\end{proposition}
\begin{proof}
Formule de Taylor en $y$, avec $\frac{d}{dy} v^{t} = -\frac{t}{f}\,v^{t+1}$ et
$\frac{d^2}{dy^2} v^t = \frac{t(t+1)}{f^2}\, v^{t+2}$.
\end{proof}

\begin{example}[La note à dix ans]\label{ex:m2:government-bonds:risk}
À 4.200~\% la note de l'\cref{ex:m2:government-bonds:accrued} affiche un \omterm{def:m2:government-bonds:accrued}{prix
pied de coupon} de 100.397405 (100-12+), un \omterm{def:m2:government-bonds:accrued}{prix plein coupon} de 100.870911, une \omterm{def:m2:government-bonds:duration}{duration de Macaulay}
de 8.142 ans, une \omterm{def:m2:government-bonds:duration}{duration modifiée} de 7.975, un \omterm{def:m2:government-bonds:duration}{DV01} de 804~USD par million et une
\omterm{def:m2:government-bonds:duration}{convexité} de 75.8. Une hausse de 100 points de base coûte 7.675 points~; la duration seule
prédit 8.044, la duration et la \omterm{def:m2:government-bonds:duration}{convexité} ensemble 7.662
(\cref{fig:m2:government-bonds:priceyield}).
\end{example}

\begin{omfigure}
\begin{tikzpicture}
\begin{axis}[width=11cm,height=5.2cm,
  xlabel={rendement (\%)},ylabel={prix plein coupon},
  tick label style={font=\small},label style={font=\small},
  xmin=1.7,xmax=6.7,ymin=80,ymax=126,grid=major,grid style={gray!25},
  legend columns=3,legend style={font=\footnotesize,at={(0.5,-0.3)},anchor=north,draw=none,fill=none,/tikz/every even column/.append style={column sep=8pt}},
  table/col sep=comma]
\addplot[omDef,very thick] table[x=yield,y=price]{figdata/markets-2/03-government-bonds/priceyield.csv};
\addplot[omThm,thick,dashed] table[x=yield,y=duration]{figdata/markets-2/03-government-bonds/priceyield.csv};
\addplot[omProp,thick,dotted] table[x=yield,y=convexity]{figdata/markets-2/03-government-bonds/priceyield.csv};
\legend{prix,duration (tangente),duration et convexité}
\end{axis}
\end{tikzpicture}

\omcaption{Prix en fonction du rendement pour la note à dix ans de
l'\cref{ex:m2:government-bonds:risk}, à deux points et demi de pourcentage de part et d'autre
de 4.20~\%. La tangente (duration) reste sous la courbe des deux côtés~: une
obligation gagne davantage quand les rendements baissent qu'elle ne perd quand ils montent
d'autant, ici de 2.2 à 2.6 points aux bords. L'approximation du second ordre
reste proche de la courbe sur l'ensemble de l'intervalle. Données~: le tutoriel du chapitre.}\label{fig:m2:government-bonds:priceyield}
\end{omfigure}

\begin{method}[Couvrir une obligation par une autre]\label{met:m2:government-bonds:hedge}
Pour couvrir une valeur nominale $N_1$ de l'obligation 1 contre les mouvements parallèles de l'obligation 2, vendre $N_2
= N_1\,\mathrm{DV01}_1/\mathrm{DV01}_2$ de valeur nominale. La couverture laisse subsister~:
la différence de \omterm{def:m2:government-bonds:duration}{convexité}, le risque que les deux rendements ne bougent pas ensemble
(risque de courbe), et la dérive des deux \omterm{def:m2:government-bonds:duration}{DV01} au fil du temps et des variations de rendement,
ce qui impose un rééquilibrage.
\end{method}

\section{Taux zéro-coupon, taux au pair et STRIPS}

\begin{definition}[Taux zéro-coupon, taux au pair, STRIPS]\label{def:m2:government-bonds:zero}
Le \emph{taux zéro-coupon}\index{taux zéro-coupon} $z(T)$ pour l'échéance $T$
est le rendement d'un paiement unique en $T$~: $P(0,T) = (1 + z/f)^{-fT}$. Le
\emph{taux au pair}\index{taux au pair} pour $T$ est le \omterm{def:m2:government-bonds:coupon}{coupon} qui ferait évaluer une
obligation arrivant à échéance en $T$ à 100. Les \emph{STRIPS}\index{STRIPS} sont les titres
\omterm{def:m2:government-bonds:coupon}{zéro-coupon} du Trésor américain, créés en séparant un billet ou une
obligation éligible en ses paiements individuels de \omterm{def:m2:government-bonds:coupon}{coupon} et de principal, chacun étant ensuite
négocié séparément~; un ensemble complet peut être recombiné pour reconstituer
le titre d'origine.
\end{definition}

\begin{proposition}[Taux au pair à partir des facteurs d'actualisation, et réciproquement]\label{prop:m2:government-bonds:par}
Avec des facteurs d'actualisation $P_k = P(0, k/f)$,
\[
\text{pair}(T_n) \;=\; f\,\frac{1 - P_n}{\sum_{k=1}^n P_k},
\qquad
P_n \;=\; \frac{1 - \frac{\text{pair}_n}{f}\sum_{k<n} P_k}{1 + \frac{\text{pair}_n}{f}} .
\]
La seconde formule \emph{bootstrappe} les facteurs d'actualisation une échéance à la fois
à partir des \omterm{def:m2:government-bonds:zero}{taux au pair}.
\end{proposition}
\begin{proof}
Une obligation de \omterm{def:m2:government-bonds:coupon}{coupon} $c$ au \omterm{def:m2:government-bonds:coupon}{pair} vérifie $100 = \frac cf \sum_{k \le n} P_k +
100 P_n$~; résoudre en $c/100$ ou, sachant $c/100 = \text{pair}_n$ et $P_1, \dots,
P_{n-1}$, en $P_n$.
\end{proof}

\begin{omfigure}
\begin{tikzpicture}
\begin{axis}[width=11cm,height=4.8cm,
  xlabel={échéance (années)},ylabel={taux (\%)},
  tick label style={font=\small},label style={font=\small},
  xmin=0,xmax=10.5,ymin=3.88,ymax=4.26,grid=major,grid style={gray!25},
  legend columns=2,legend style={font=\footnotesize,at={(0.5,-0.3)},anchor=north,draw=none,fill=none,/tikz/every even column/.append style={column sep=8pt}},
  table/col sep=comma]
\addplot[omDef,very thick,mark=*,mark size=1.2pt] table[x=years,y=par]{figdata/markets-2/03-government-bonds/zero.csv};
\addplot[omThm,very thick,dashed] table[x=years,y=zero]{figdata/markets-2/03-government-bonds/zero.csv};
\legend{taux au pair,taux zéro-coupon (bootstrappés)}
\end{axis}
\end{tikzpicture}

\omcaption{Une courbe des \omterm{def:m2:government-bonds:zero}{taux au pair} illustrative et la courbe \omterm{def:m2:government-bonds:coupon}{zéro-coupon} bootstrappée à partir de celle-ci
(\cref{prop:m2:government-bonds:par}). Lorsque la courbe des \omterm{def:m2:government-bonds:zero}{taux au pair} monte, la courbe \omterm{def:m2:government-bonds:coupon}{zéro-coupon}
se situe au-dessus d'elle, car le rendement d'une obligation à \omterm{def:m2:government-bonds:coupon}{coupons} moyenne les \omterm{def:m2:government-bonds:zero}{taux zéro-coupon} de
tous ses flux, la plupart étant plus proches et aux taux plus faibles~; là où la courbe des \omterm{def:m2:government-bonds:zero}{taux au pair} fléchit,
entre un et deux ans, la courbe \omterm{def:m2:government-bonds:coupon}{zéro-coupon} passe sous elle. Données~: celles
du tutoriel du chapitre.}\label{fig:m2:government-bonds:zero}
\end{omfigure}

\begin{omfigure}
\begin{tikzpicture}
\begin{axis}[width=11cm,height=4.8cm,
  xlabel={jours après le 15 août 2026},ylabel={prix},
  tick label style={font=\small},label style={font=\small},
  xmin=0,xmax=364,ymin=100.2,ymax=102.8,grid=major,grid style={gray!25},
  legend columns=2,legend style={font=\footnotesize,at={(0.5,-0.3)},anchor=north,draw=none,fill=none,/tikz/every even column/.append style={column sep=8pt}},
  table/col sep=comma]
\addplot[omThm,very thick] table[x=day,y=dirty]{figdata/markets-2/03-government-bonds/cleandirty.csv};
\addplot[omDef,very thick] table[x=day,y=clean]{figdata/markets-2/03-government-bonds/cleandirty.csv};
\legend{prix plein coupon,prix pied de coupon}
\end{axis}
\end{tikzpicture}

\omcaption{Le billet à dix ans sur un an, à rendement constant de 4.20~\%, calculé chaque jour.
Le \omterm{def:m2:government-bonds:accrued}{prix plein coupon} croît à mesure que le \omterm{def:m2:government-bonds:accrued}{coupon couru} s'accumule et chute du montant du \omterm{def:m2:government-bonds:coupon}{coupon} lorsque celui-ci
est versé (15 février, jour 184)~; le \omterm{def:m2:government-bonds:accrued}{prix pied de coupon} ne bouge que par sa lente
convergence vers le \omterm{def:m2:government-bonds:coupon}{pair}. Données~: celles du tutoriel du chapitre.}\label{fig:m2:government-bonds:cleandirty}
\end{omfigure}

\section{Tutoriel~: une obligation sur un calendrier réel}\label{tut:m2:government-bonds:bond}

\begin{tutorial}
\textbf{Objectif.} Calculer le prix du billet à dix ans à partir de son rendement au 25 septembre 2026,
vérifier son \omterm{def:m2:government-bonds:duration}{DV01} par rapport à une revalorisation avec choc, et bootstrapper une courbe \omterm{def:m2:government-bonds:coupon}{zéro-coupon}.
\textbf{État final~:} les résultats de l'\cref{ex:m2:government-bonds:risk} et
les \cref{fig:m2:government-bonds:priceyield,fig:m2:government-bonds:zero}.
\begin{tutsteps}
\item \textbf{L'échéancier et le \omterm{def:m2:government-bonds:accrued}{coupon couru}}, généré à rebours à partir de
l'échéance, avec la règle de fin de mois.
\omcode{code/firm/bond/firm_bond.py}{51}{71}{Dates de coupon, fraction de la période restant à courir, coupon couru.}\label{lst:m2:government-bonds:schedule}
Il faut obtenir 0.473505 pour le billet au 25 septembre 2026.
\item \textbf{Le prix} dans les deux conventions de la \cref{rem:m2:government-bonds:conventions}.
\omcode{code/firm/bond/firm_bond.py}{73}{78}{Prix plein coupon à partir du rendement~: convention street et convention du Trésor.}\label{lst:m2:government-bonds:dirty}
Il faut vérifier les exemples de la réglementation~: 99.057893 et 99.730918.
\item \textbf{Risque}, dans le langage système de la firme. Le jumeau C++20
calcule la duration, le \omterm{def:m2:government-bonds:duration}{DV01} et la \omterm{def:m2:government-bonds:duration}{convexité} en un seul passage sur les flux.
\omcode{code/firm/bond/cpp/firm_bond.hpp}{94}{108}{Duration, DV01 et convexité (C++20).}\label{lst:m2:government-bonds:risk}
Le \omterm{def:m2:government-bonds:duration}{DV01} doit égaler la moitié de la différence de prix entre les rendements
décalés de 1 point de base de part et d'autre, à six chiffres significatifs près.
\item \textbf{La courbe \omterm{def:m2:government-bonds:coupon}{zéro-coupon}.}
\omcode{code/firm/bond/firm_bond.py}{116}{122}{Bootstrap des facteurs d'actualisation à partir des taux au pair.}\label{lst:m2:government-bonds:bootstrap}
\texttt{fig\_bond.py} écrit les trois graphiques du chapitre.
\end{tutsteps}
\textbf{Pour aller plus loin.} Calculer le prix du même billet au 13 février 2027 et au
16 février 2027, et comparer les \omterm{def:m2:government-bonds:accrued}{prix pied de coupon} et plein \omterm{def:m2:government-bonds:coupon}{coupon} de part et d'autre du versement du \omterm{def:m2:government-bonds:coupon}{coupon}~; puis
fixer le \omterm{def:m2:government-bonds:coupon}{coupon} à zéro et vérifier que la \omterm{def:m2:government-bonds:duration}{duration de Macaulay} est égale à la
durée de vie restante.
\end{tutorial}

\section{Construction~: la bibliothèque des obligations}\label{bld:m2:government-bonds:bond}

\begin{build}
\bfield{Rôle} Le premier élément de la bibliothèque de valorisation de la firme miniature
(étendue dans les livres~5 et~6)~: chaque composante de taux, la base des contrats à terme
(du \cref{ch:m2:bond-futures}), l'analyseur d'adjudications
(du \cref{ch:m2:the-treasury-market}), les mesures de spread de la partie crédit
(du \cref{ch:m2:corporate-bonds}), valorise les obligations par son intermédiaire. Python pour la recherche,
C++20 pour la chaîne de trading, Rust comme jumeau.
\bfield{Interface} \texttt{Bond(coupon, maturity, freq)} avec
\texttt{coupon\_dates(settle)}, \texttt{accrued(settle)},
\texttt{dirty\_price(y, settle, treasury)}, \texttt{clean\_price(\ldots)},
\texttt{yield\_from\_clean(clean, settle)}, \texttt{risk(y, settle)}~;
\texttt{act\_360}, \texttt{act\_365f}, \texttt{thirty\_360}~;
\texttt{bootstrap\_par(pars)}, \texttt{par\_yield(dfs)}. L'en-tête C++ et la
crate Rust exposent les mêmes noms.
\bfield{Règles} Dates de \omterm{def:m2:government-bonds:coupon}{coupon} calculées à rebours depuis l'échéance, fin de mois conservée~;
\omterm{def:m2:government-bonds:accrued}{coupon couru} actual/actual~; la convention street par défaut et la
convention du Trésor sur demande~; rendements décimaux, capitalisés à
la fréquence du \omterm{def:m2:government-bonds:coupon}{coupon}.
\bfield{Tests d'acceptation} \texttt{code/firm/bond/tests/},
\texttt{cpp/firm\_bond\_test.cpp}, \texttt{rust/}~: les exemples traités de la
réglementation (99.057893~; \omterm{def:m2:government-bonds:accrued}{coupon couru} 0.367403 et 99.730918)~; un échéancier de
fin de mois en année bissextile~; l'aller-retour prix-rendement~; le \omterm{def:m2:government-bonds:duration}{DV01} et la
\omterm{def:m2:government-bonds:duration}{convexité} face à des chocs~; la duration d'une obligation \omterm{def:m2:government-bonds:coupon}{zéro-coupon} est égale à sa durée de vie~; les facteurs
bootstrappés revalorisent leurs \omterm{def:m2:government-bonds:zero}{taux au pair}.
\bfield{Pour aller plus loin} Premiers \omterm{def:m2:government-bonds:coupon}{coupons} courts et longs (les autres
exemples de la réglementation)~; périodes ex-dividende comme pour les gilts~; la variante européenne 30/360~; un
calendrier des jours ouvrés pour les dates de paiement.
\end{build}

\begin{omsources}
\item 31 CFR Part 356, Appendix B, \emph{Formulas and Tables}, sections I et
II. US Treasury, \emph{STRIPS} (TreasuryDirect)~; \emph{Monthly Statement of
the Public Debt}, août 2026.
\item UK Debt Management Office, \emph{Formulae for Calculating Gilt Prices from Yields}. Deutsche Finanzagentur, \emph{Federal Bonds}.
\item P.\ Ritchken, \emph{Fixed Income}, chapitre 2, «~Treasury Securities~».
\item F.\,J.\ Fabozzi (dir.), \emph{The Handbook of Fixed Income Securities},
McGraw-Hill, pour les conventions des autres marchés.
\end{omsources}

\section{Exercices}

\begin{exercise}[$\star$]\label{exo:m2:government-bonds:1}
Un billet à 4.25~\% à \omterm{def:m2:government-bonds:coupon}{coupon} semestriel verse un \omterm{def:m2:government-bonds:coupon}{coupon} le 15 août. Donner son \omterm{def:m2:government-bonds:accrued}{coupon couru}
pour 100 de valeur nominale, pour un règlement-livraison le 25 septembre, la période étant de 184 jours.
\end{exercise}

\begin{exercise}[$\star$]\label{exo:m2:government-bonds:2}
Convertir 99-16+ en prix décimal, et 101.296875 en cotation en 32e.
\end{exercise}

\begin{exercise}[$\star$]\label{exo:m2:government-bonds:3}
Donner le \omterm{def:m2:government-bonds:accrued}{prix pied de coupon}, à une date de \omterm{def:m2:government-bonds:coupon}{coupon}, d'une obligation à trois ans à 4~\% à \omterm{def:m2:government-bonds:coupon}{coupon} semestriel,
pour un rendement de 4~\%, puis de 5~\%.
\end{exercise}

\begin{exercise}[$\star\star$]\label{exo:m2:government-bonds:4}
Pour une obligation à deux ans à 4~\% à \omterm{def:m2:government-bonds:coupon}{coupon} semestriel, pour un rendement de 4~\% à une date de \omterm{def:m2:government-bonds:coupon}{coupon}, donner
la \omterm{def:m2:government-bonds:duration}{duration de Macaulay}, la \omterm{def:m2:government-bonds:duration}{duration modifiée} et le \omterm{def:m2:government-bonds:duration}{DV01} par million.
\end{exercise}

\begin{exercise}[$\star\star$]\label{exo:m2:government-bonds:5}
Pour le billet à dix ans de l'\cref{ex:m2:government-bonds:risk}, estimer la
variation de prix pour une \emph{baisse} de rendement de 100 points de base, avec la seule duration,
puis avec la \omterm{def:m2:government-bonds:duration}{convexité}, et comparer avec une valorisation complète.
\end{exercise}

\begin{exercise}[$\star\star$]\label{exo:m2:government-bonds:6}
Une obligation à 5~\% à \omterm{def:m2:government-bonds:coupon}{coupon} semestriel verse ses \omterm{def:m2:government-bonds:coupon}{coupons} le 15 février et le 15 août. Calculer son \omterm{def:m2:government-bonds:accrued}{coupon
couru} pour un règlement-livraison le 31 mars 2026, selon les conventions actual/actual puis 30/360.
Quel investisseur, sur quel marché, se soucierait de cet écart~?
\end{exercise}

\begin{exercise}[$\star\star\star$]\label{exo:m2:government-bonds:7}
\emph{Codage.} Avec \texttt{bootstrap\_par}, bootstrapper la courbe des \omterm{def:m2:government-bonds:zero}{taux au pair} de
la \cref{fig:m2:government-bonds:zero}, puis donner le facteur d'actualisation et
le \omterm{def:m2:government-bonds:zero}{taux zéro-coupon} à dix ans. Expliquer le signe de l'écart entre \omterm{def:m2:government-bonds:zero}{taux zéro-coupon} et \omterm{def:m2:government-bonds:zero}{taux au pair} à dix ans.
\end{exercise}

\begin{exercise}[$\star\star\star$]\label{exo:m2:government-bonds:8}
\emph{Trouver l'erreur.} «~Une nouvelle obligation à trente ans a une duration de trente
ans, donc une hausse d'un point de base des rendements coûte 0.30~\% de son prix.~»
Corriger cette affirmation pour une obligation à 4.5~\% à \omterm{def:m2:government-bonds:coupon}{coupon} semestriel émise au \omterm{def:m2:government-bonds:coupon}{pair}.
\end{exercise}

\section{Problème~: le desk à deux obligations}

\begin{problem}[{Problème du week-end --- une position de courbe neutre en DV01}]\label{pb:m2:government-bonds:1}
Le 25 septembre 2026, un desk achète 100 millions de dollars de valeur nominale d'un billet à cinq ans (4~\%
\omterm{def:m2:government-bonds:coupon}{coupon}, échéance le 30 septembre 2031, rendement de 3.95~\%) et le couvre avec le billet à dix ans
de l'\cref{ex:m2:government-bonds:risk} (4.25~\%, 15 août 2036, rendement de
4.20~\%).

\medskip\noindent\textbf{Partie I --- Prix.}
\begin{enumerate}
\item Donner les dates de \omterm{def:m2:government-bonds:coupon}{coupon} précédente et suivante du billet à cinq ans. Pourquoi son
prochain \omterm{def:m2:government-bonds:coupon}{coupon} est-il si proche~?
\item Donner son \omterm{def:m2:government-bonds:accrued}{coupon couru} et son \omterm{def:m2:government-bonds:accrued}{prix pied de coupon}, en décimal et en 32e.
\item Donner son \omterm{def:m2:government-bonds:accrued}{prix plein coupon} pour 100 millions de dollars de valeur nominale.
\item Donner le \omterm{def:m2:government-bonds:duration}{DV01} par million de chacun des deux billets.
\item Donner la \omterm{def:m2:government-bonds:duration}{convexité} de chacun des deux billets.
\end{enumerate}

\medskip\noindent\textbf{Partie II --- La couverture.}
\begin{enumerate}[resume]
\item Quelle valeur nominale du billet à dix ans le desk doit-il vendre pour être neutre en \omterm{def:m2:government-bonds:duration}{DV01}~?
\item Quel est le flux de trésorerie net de la position (\omterm{def:m2:government-bonds:accrued}{prix plein coupon} de la jambe longue moins
celui de la jambe courte)~?
\item Revaloriser les deux billets avec les deux rendements majorés de 1 point de base. Quel
est le P\&L~?
\item Revaloriser avec les deux rendements majorés de 25 points de base, puis minorés de 25 points de base. Quel
est le P\&L dans chaque cas~?
\item Pourquoi est-il négatif dans les deux sens~?
\end{enumerate}

\medskip\noindent\textbf{Partie III --- La courbe.}
\begin{enumerate}[resume]
\item Le rendement à cinq ans monte de 10 points de base, celui à dix ans ne bouge pas.
Quel est le P\&L~?
\item Le rendement à dix ans monte de 10 points de base, celui à cinq ans ne bouge pas.
Quel est le P\&L~?
\item Quel mouvement de courbe profite à la position~? Le nommer.
\item La duration de la position est-elle nulle~? Son \omterm{def:m2:government-bonds:duration}{DV01}~?
\item Quatre semaines plus tard, aucun des deux rendements n'a bougé. La position est-elle encore
neutre en \omterm{def:m2:government-bonds:duration}{DV01}~? Qu'est-ce qui a changé~?
\end{enumerate}

\medskip\noindent\textbf{Partie IV --- Jugement.}
\begin{enumerate}[resume]
\item Pourquoi couvrir le \omterm{def:m2:government-bonds:duration}{DV01} plutôt que la valeur nominale ou la valeur de marché~?
\item Nommer deux risques que la couverture laisse subsister.
\item Pourquoi le desk cote-t-il les deux billets en 32e mais calcule-t-il en décimal~?
\item Énoncer le \emph{résultat nommé}~: le ratio de couverture et le P\&L de la
position couverte pour un mouvement parallèle de 25 points de base dans un sens ou dans l'autre.
\item En une phrase~: sur quoi parie une position neutre en \omterm{def:m2:government-bonds:duration}{DV01}~?
\end{enumerate}
\end{problem}

\section{Questions d'entretien}

\begin{interviewq}[$\star$ \iqroles{trader, banque}]\label{iq:m2:government-bonds:1}
Quelle est la différence entre le \omterm{def:m2:government-bonds:accrued}{prix pied de coupon} et le \omterm{def:m2:government-bonds:accrued}{prix plein coupon} d'une
obligation, et pourquoi le marché cote-t-il le premier~?
\end{interviewq}

\begin{interviewq}[$\star$ \iqroles{trader, chercheur}]\label{iq:m2:government-bonds:2}
Pourquoi la relation prix-rendement est-elle convexe~? La \omterm{def:m2:government-bonds:duration}{convexité} est-elle favorable au détenteur~?
\end{interviewq}

\begin{interviewq}[$\star\star$ \iqroles{trader, chercheur, banque}]\label{iq:m2:government-bonds:3}
Une obligation \omterm{def:m2:government-bonds:coupon}{zéro-coupon} à dix ans et une obligation à \omterm{def:m2:government-bonds:coupon}{coupon} de 5~\% à dix ans~: laquelle a la
plus grande duration~? Le plus grand \omterm{def:m2:government-bonds:duration}{DV01} pour 100 de valeur nominale~?
\end{interviewq}

\begin{interviewq}[$\star\star$ \iqroles{trader}]\label{iq:m2:government-bonds:4}
Vous êtes acheteur sur les cinq ans et voulez couvrir avec les dix ans. Combien
vendez-vous, et quels risques subsistent~?
\end{interviewq}

\begin{interviewq}[$\star\star$ \iqroles{développeur}]\label{iq:m2:government-bonds:5}
Vous devez implémenter le calcul du \omterm{def:m2:government-bonds:accrued}{coupon couru} pour les obligations d'État dans trois
marchés. Que couvriront vos tests~?
\end{interviewq}

\begin{interviewq}[$\star\star\star$ \iqroles{chercheur, trader}]\label{iq:m2:government-bonds:6}
Quelle est la \omterm{def:m2:government-bonds:duration}{duration modifiée} d'une obligation perpétuelle versant un \omterm{def:m2:government-bonds:coupon}{coupon} fixe une
fois par an, à un rendement de 4~\%~? Et sa \omterm{def:m2:government-bonds:duration}{duration de Macaulay}~?
\end{interviewq}
